\documentclass[11pt]{article}

\usepackage[margin=1in]{geometry}
\usepackage{amsmath}
\usepackage{booktabs}
\usepackage{hyperref}
\usepackage{enumitem}

\title{Methodology of the xGelo Forecasting Model}
\author{xGelo Project}
\date{June 20, 2026}

\begin{document}
\maketitle

\begin{abstract}
xGelo is an open-data football forecasting system for international fixtures,
with a particular focus on UEFA World Cup Qualifiers and the 2026 World Cup. The
implemented model combines a chronological Elo rating system, a shot-level
expected-goals model, leakage-safe team form and squad-strength features,
negative-binomial goal models, and Monte Carlo simulation. This paper summarizes
the methodology implemented in the repository, including the distinction between
the baseline xG+Elo pipeline and the current 2026 World Cup dashboard deployment.
The baseline pipeline retains the original xG and rolling-form components, while
the current dashboard uses Elo, weighted historical goal ability, and
Transfermarkt player-pool features as active goal-model predictors because the
available international rolling xG/form coverage is insufficient.
\end{abstract}

\section{Model Overview}

xGelo is designed as a reproducible analytical pipeline rather than as a
black-box prediction service. The pipeline is defined in \path{_targets.R} and
uses R targets to materialize data, fitted models, validation outputs, and
dashboard artifacts. The model has five conceptual layers:

\begin{enumerate}[leftmargin=*]
  \item data ingestion from open or local-only sources;
  \item shot-level expected-goals estimation from StatsBomb event data;
  \item chronological Elo ratings from historical international results;
  \item match-level feature construction with strict pre-match lookups; and
  \item negative-binomial score simulation for win, draw, loss, scoreline, and
  tournament probabilities.
\end{enumerate}

The implementation has two closely related forecast routes. The baseline route
uses Elo differences and candidate rolling xG/form features. The current
WC2026 dashboard route uses a hybrid feature set: Elo difference, weighted
historical attack and defense ability, and leakage-safe Transfermarkt
player-pool valuation features. The overall approach sits in the long
tradition of football score modeling and tournament simulation, from early
attack-defense Poisson models through modern Elo, ranking, and hybrid
tournament forecast systems
\cite{maher1982,dixoncoles1997,karlisntzoufras2003,gilchmuller2018,groll2018rf}.
This hybrid design is conceptually close to the
tournament forecasting work of Groll, Ley, Schauberger, Van Eetvelde, Zeileis,
and coauthors, who combine estimated team abilities, player or squad
information, covariates, score prediction, and repeated tournament simulation
\cite{groll2019wwc,groll2021euro2020,groll2024euro2024}. xGelo adapts that
idea to a deliberately simpler open-data setting: model-derived strength
signals are used as covariates in goal-count models, but the active score model
is a pair of interpretable negative-binomial GLMs rather than a random forest,
xgboost model, or ensemble. The shot-level xG model remains trained and audited
in the pipeline, but its derived rolling international form variables are
inactive in the current dashboard artifact because the generated
\path{data/processed/xg_feature_usage_audit.csv} records zero variance and no
forecast-team rolling-form coverage for those predictors.

\section{Data Sources and Temporal Discipline}

The primary historical match source is the martj42 international results data
set, which provides match dates, teams, venues, neutral-site flags, scores, and
tournaments from 1872 onward. In the current local artifacts,
\path{data/processed/elo_matches.csv} contains 49,477 rows from 1872-11-30 to
2026-06-27, and \path{data/processed/elo_ratings.csv} contains 198,084 rating
history rows through 2026-06-19. StatsBomb Open Data supplies event-level shot
coordinates for the xG model. Transfermarkt-derived player-pool features are
optional and local-only: the raw DuckDB snapshot is not redistributed, while
derived squad-strength features are built from valuations dated strictly before
the target match or benchmark cutoff.

Temporal leakage prevention is a central part of the implementation. Feature
tables are constructed using ``latest before'' lookup functions for Elo,
rolling form, squad strength, and goal ability. A feature row includes a
\texttt{feature\_source\_date}, and \texttt{assert\_no\_feature\_leakage()}
fails if that date is on or after the fixture date. For fixed-cutoff forecasts,
such as the WC2026 dashboard, all features are sourced before the cutoff. The
current WC2026 feature artifact uses cutoff 2026-06-19 and contains 40 future
group-fixture rows plus 2,256 ordered knockout candidate rows.

\section{Shot-Level Expected Goals}

The xG component estimates the probability that a shot becomes a goal. For each
StatsBomb shot, the implementation extracts:

\begin{itemize}[leftmargin=*]
  \item distance to the center of the goal;
  \item shot angle subtended by the two posts;
  \item whether the shot is a header;
  \item whether the shot comes from regular open play; and
  \item competition identity.
\end{itemize}

Penalties are excluded. StatsBomb coordinates are treated as a 120-by-80 yard
pitch with the goal centered at $(120,40)$ and posts at $y=36.34$ and
$y=43.66$. For a shot at $(x,y)$, distance is
\[
d = \sqrt{(120 - x)^2 + (40 - y)^2}.
\]
The shot angle is computed with the law of cosines. Let $a$ and $b$ be the
distances from the shot location to the two posts and $c$ be the goal width.
Then
\[
\theta = \cos^{-1}\left(\frac{a^2 + b^2 - c^2}{2ab}\right),
\]
with the cosine argument clamped to $[-1,1]$ for numerical stability.

The fitted xG model is a logistic regression implemented with tidymodels. Natural
spline bases with four degrees of freedom are applied to distance and angle, and
categorical predictors are dummy encoded after zero-variance filtering:
\[
\Pr(\mathrm{goal}_i = 1)
  = \mathrm{logit}^{-1}\left(\beta_0
  + f_d(d_i) + f_\theta(\theta_i) + \beta^\top z_i\right),
\]
where $z_i$ contains the header, open-play, and competition indicators. The
current backtest output reports AUC 0.7905, accuracy 0.8462, and Brier score
0.1188 in \path{outputs/model_performance/xg_backtest.csv}.

\section{Chronological Elo Ratings}

The Elo system is a custom implementation for international football. Teams
start at 1500. Matches are processed in chronological order, and each team has a
pre-match and post-match rating history entry. The rating update follows the
general Elo framework \cite{elo1978}, which has been adapted widely for
football ranking and prediction \cite{lasek2013,gilchmuller2018}. More general
dynamic team-rating approaches, such as Knorr-Held's dynamic sports-team rating
model, motivate the broader idea that team strength should evolve through time
rather than be estimated as a static season average \cite{knorrheld2000}. The
expected result for team $A$ against team $B$ is
\[
E_A = \frac{1}{1 + 10^{(R_B - R_A)/400}}.
\]
For non-neutral matches, a 60-point home advantage is added to the home team's
pre-match rating before the expected result is computed. The rating update is
\[
R_A' = R_A + K_A(S_A - E_A),
\]
where $S_A$ is 1 for a win, 0.5 for a draw, and 0 for a loss. The implementation
uses $K=20$ for teams with at least 15 matches across the previous and current
year counters, and $K=40$ otherwise. Ratings also decay with inactivity:
\[
R(t) = R_0 \cdot 0.995^{\Delta t / 365},
\]
where $\Delta t$ is the number of days since the team's last match.

The validation artifact reports an Elo home-win discrimination AUC of 0.7916.
In later forecast layers, Elo enters as \texttt{elo\_diff}, the pre-match home
rating minus the pre-match away rating after applying venue adjustment where
appropriate. This mirrors the football forecasting literature in which Elo or
ranking-derived team strength is used either directly as a prediction rule or
as a covariate in score-count models and tournament simulations
\cite{lasek2013,gilchmuller2018,groll2018rf}.

\section{Match-Level Feature Construction}

The feature layer converts source data into one row per fixture. In the baseline
route, candidate predictors include differences in rolling xG for, xG against,
xG difference, shots, and a composite form index. Rolling form is computed with
an exponentially weighted moving average:
\[
\mathrm{EWMA}_t = \alpha x_t + (1-\alpha)\mathrm{EWMA}_{t-1},
\qquad
\alpha = \frac{2}{\mathrm{span}+1},
\]
with a default span of 12 matches. Importantly, post-match EWMAs are lagged so
that each row contains only pre-match form.

The current hybrid route expands the match row with two additional families of
features. The motivation follows the hybrid tournament-forecasting literature:
instead of relying on one rating or one set of raw covariates, the score model
is supplied with several partially independent strength summaries. In the
Zeileis-associated EURO and World Cup work, these summaries include historical
team-ability estimates, bookmaker-consensus ability estimates, player ratings,
market values, team structure, and other team covariates
\cite{leitner2010,groll2015,groll2019wwc,groll2021euro2020,groll2024euro2024}.
xGelo implements a smaller version of this idea that can be reproduced from the
local data available in the project.

First, weighted historical goal ability summarizes attack and defense using
prior international scores. Let $w_m$ be the product of a recency weight and a
tournament-importance weight for historical match $m$. The recency component
uses a 730-day half-life,
\[
w_m^{\mathrm{time}} = \exp\left(-\log(2)
  \frac{\mathrm{age}_m}{730}\right),
\]
and the tournament component gives major tournaments, qualifiers, and Nations
League matches more influence than friendlies. For each team, weighted goals
for and weighted goals against are divided by their corresponding global
weighted rates. Thus an attack value above 1 indicates scoring above the global
baseline after weighting, while a defense value above 1 indicates conceding
above the global baseline. Fixture-level predictors are then formed as
\texttt{attack\_ability\_diff} and \texttt{defense\_ability\_diff}, measured
from the listed home-team perspective.

Second, Transfermarkt player-pool strength approximates squad quality and depth.
The implementation intentionally avoids current-only leakage: raw player profile
fields that are not historically dated are dropped for historical benchmarks,
and player valuations are included only when their valuation date precedes the
match or forecast cutoff. The resulting features are not a single market-value
number. They include log total squad value, top-11, top-15, and top-23 value,
median player value, value concentration, positional value aggregates for
goalkeepers, defense, midfield, and attack, age/value composition, and 6- and
12-month momentum. Each is converted to a home-minus-away difference before
entering the model.

The fitted hybrid models retain 26 predictors: \texttt{elo\_diff}, attack and
defense ability differences, and the Transfermarkt value, structure, position,
age, and momentum differences listed by \path{R/forecast/poisson.R}. The
home-goal and away-goal models use the same home-perspective feature convention;
for neutral fixtures, prediction averages over the original and reversed
perspectives to avoid giving the first-listed team an artificial home-model
intercept advantage. Candidate rolling xG/form features are present in the code
but not active in the current fitted hybrid models because their current
international feature table values have zero variance and no WC2026 forecast
team coverage.

\section{Goal Models}

Statistical analysis of football scores predates modern machine-learning
tournament forecasts. Early work on football score distributions and chance
includes Moroney, Reep and Benjamin, and Hill
\cite{moroney1956,reepbenjamin1968,hill1974}. The modern score-model lineage
then centers on count models for home and away goals, especially Poisson and
related bivariate or adjusted count models.

The forecast engine fits separate goal-count models for the listed home and away
teams. The preferred model is a negative-binomial generalized linear model
fitted with \texttt{MASS::glm.nb}; if the negative-binomial fit fails, the code
falls back to a Poisson GLM. This score-count framing follows the standard
football modeling lineage, including independent attack-defense Poisson models,
low-score dependence corrections, bivariate Poisson models, and Skellam
goal-difference models
\cite{maher1982,lee1997,dixoncoles1997,karlisntzoufras2003,karlisntzoufras2009}.
Recent work continues to study dependence and distributional assumptions in
football score models \cite{petretta2021,michels2023}. For side
$s \in \{\mathrm{home}, \mathrm{away}\}$,
\[
Y_s \sim \mathrm{NegBin}(\mu_s, \theta_s), \qquad
\log(\mu_s) = \beta_{0s} + \beta_s^\top x.
\]
The negative-binomial form is used because football scores are low-count but
overdispersed relative to a strict Poisson assumption.

The saved baseline models use only \texttt{elo\_diff}. Their fitted dispersion
parameters are $\theta=6.083$ for home goals and $\theta=3.010$ for away goals.
The saved hybrid models use the retained predictor set described above, with
$\theta=6.114$ for home goals and $\theta=4.945$ for away goals. The hybrid
training table currently contains 49,437 scored international matches through
2026-06-19, with 24,229 home wins, 11,243 draws, and 13,965 away wins.

\section{Monte Carlo Match and Tournament Simulation}

For each forecast fixture, the implementation loads the relevant home and away
goal models, constructs the fixture feature row, predicts mean goals, and draws
simulated scores from the fitted negative-binomial distributions:
\[
Y_\mathrm{home}^{(m)} \sim \mathrm{NegBin}(\mu_\mathrm{home}, \theta_\mathrm{home}),
\qquad
Y_\mathrm{away}^{(m)} \sim \mathrm{NegBin}(\mu_\mathrm{away}, \theta_\mathrm{away}).
\]
The empirical frequencies across $M$ simulations define home-win, draw, and
away-win probabilities:
\[
\hat{p}_{H} = \frac{1}{M}\sum_{m=1}^{M} \mathbf{1}
\{Y_\mathrm{home}^{(m)} > Y_\mathrm{away}^{(m)}\},
\]
with analogous definitions for draws and away wins. The same score draws also
produce expected goals, modal scorelines, top exact-score probabilities,
over/under 2.5 goals, and both-teams-to-score probabilities.

Neutral matches require special handling because separate home and away model
intercepts otherwise create an arbitrary boost for the first-listed team. For
neutral fixtures using home-perspective models, xGelo averages the appropriate
home-model and away-model predictions under the original and reversed team
perspectives.

The tournament simulator composes match forecasts into group and knockout paths.
Group-stage simulations sample exact scores, award points, and rank groups using
World Cup-style head-to-head tiebreakers before goal difference, goals for, and
a seeded residual tiebreaker. Knockout matches use the 90-minute score model;
if regulation is drawn, advancement is allocated through an Elo-based tiebreak
share representing the combined extra-time and penalties path. The dashboard can
run a fast 100,000-tournament simulator by precomputing fixture distributions
and ordered knockout advancement matrices. Repeated tournament simulation is
common in football forecasting because nonlinear tournament rules, group draws,
tiebreakers, and knockout brackets make closed-form stage probabilities
impractical \cite{kuonen1996,gilchmuller2018,groll2018rf,groll2021euro2020}.

\section{Validation and Limitations}

The implemented validation checks combine statistical metrics, unit tests, and
pipeline assertions. The xG model exceeds its target AUC threshold of 0.75. The
Elo model's validation AUC is 0.7916. Forecast calibration artifacts report a
draw-probability Brier score near 0.214 and predicted draw frequency near the
target WCQ-UEFA rate of 28 percent in the v1 project state.

Several limitations remain. The xG model is trained on available open
StatsBomb data rather than the target WCQ-UEFA competition, and the current
international rolling xG/form coverage is too sparse to be active in the
published dashboard. Transfermarkt values are useful proxies for player-pool
strength but are not a direct measure of tactical fit, player availability, or
injury status. The knockout tiebreak model combines extra time and penalties
rather than modeling them separately. Finally, all forecasts are pre-match
forecasts; the system is not designed for in-play updating.

\section{Conclusion}

xGelo implements an interpretable open-data forecasting stack: chronological Elo
captures long-run national-team strength, shot-level xG provides a transparent
event-data component, leakage-safe feature construction prevents future
information from entering forecasts, negative-binomial goal models represent
score uncertainty, and Monte Carlo simulation translates goal intensities into
match and tournament probabilities. The present WC2026 dashboard deployment is
best understood as a hybrid international-goal model using Elo, historical goal
ability, and dated squad-value information, with xG/form features retained as
audited but currently inactive candidate predictors.

\begin{thebibliography}{99}
\bibitem{martj42}
martj42. \emph{International football results from 1872 to present}.
\url{https://github.com/martj42/international_results}

\bibitem{statsbomb}
StatsBomb. \emph{StatsBomb Open Data}.
\url{https://github.com/statsbomb/open-data}

\bibitem{transfermarkt}
dcaribou. \emph{transfermarkt-datasets}.
\url{https://github.com/dcaribou/transfermarkt-datasets}

\bibitem{elo1978}
Elo, A. E. \emph{The Rating of Chessplayers, Past and Present}. Arco, 1978.

\bibitem{moroney1956}
Moroney, M. J. \emph{Facts from Figures}. 3rd ed. Penguin, 1956.

\bibitem{reepbenjamin1968}
Reep, C. and Benjamin, B. Skill and chance in association football.
\emph{Journal of the Royal Statistical Society. Series A}, 131(4):581--585,
1968.

\bibitem{hill1974}
Hill, I. D. Association football and statistical inference.
\emph{Applied Statistics}, 23(2):203--208, 1974.

\bibitem{maher1982}
Maher, M. J. Modelling association football scores. \emph{Statistica
Neerlandica}, 36(3):109--118, 1982.

\bibitem{kuonen1996}
Kuonen, D. Statistical models for knock-out soccer tournaments. Technical
report, 1996.

\bibitem{lee1997}
Lee, A. J. Modeling scores in the Premier League: Is Manchester United really
the best? \emph{Chance}, 10(1):15--19, 1997.

\bibitem{dixoncoles1997}
Dixon, M. J. and Coles, S. G. Modelling association football scores and
inefficiencies in the football betting market. \emph{Applied Statistics},
46(2):265--280, 1997.

\bibitem{knorrheld2000}
Knorr-Held, L. Dynamic rating of sports teams. \emph{The Statistician},
49(2):261--276, 2000.

\bibitem{karlisntzoufras2003}
Karlis, D. and Ntzoufras, I. Analysis of sports data by using bivariate Poisson
models. \emph{The Statistician}, 52(3):381--393, 2003.

\bibitem{karlisntzoufras2009}
Karlis, D. and Ntzoufras, I. Bayesian modelling of football outcomes: Using the
Skellam's distribution for the goal difference. \emph{IMA Journal of Management
Mathematics}, 20(2):133--145, 2009.

\bibitem{leitner2010}
Leitner, C., Zeileis, A., and Hornik, K. Forecasting sports tournaments by
ratings of (prob)abilities: A comparison for the EURO 2008. \emph{International
Journal of Forecasting}, 26(3):471--481, 2010.

\bibitem{groll2015}
Groll, A., Schauberger, G., and Tutz, G. Prediction of major international
soccer tournaments based on team-specific regularized Poisson regression: An
application to the FIFA World Cup 2014. \emph{Statistical Modelling},
15(4):335--360, 2015.

\bibitem{lasek2013}
Lasek, J., Szl\'avik, Z., and Bhulai, S. The predictive power of ranking systems
in association football. \emph{International Journal of Applied Pattern
Recognition}, 1(1):27--46, 2013.

\bibitem{gilchmuller2018}
Gilch, L. A. and M{\"u}ller, S. \emph{On Elo Based Prediction Models for the
FIFA Worldcup 2018}. arXiv:1806.01930, 2018.
\url{https://arxiv.org/abs/1806.01930}

\bibitem{groll2018rf}
Groll, A., Ley, C., Schauberger, G., and Van Eetvelde, H. \emph{Prediction of
the FIFA World Cup 2018 -- A Random Forest Approach with an Emphasis on
Estimated Team Ability Parameters}. arXiv:1806.03208, 2018.
\url{https://arxiv.org/abs/1806.03208}

\bibitem{groll2019wwc}
Groll, A., Ley, C., Schauberger, G., Van Eetvelde, H., and Zeileis, A.
\emph{Hybrid Machine Learning Forecasts for the FIFA Women's World Cup 2019}.
arXiv:1906.01131, 2019.
\url{https://arxiv.org/abs/1906.01131}

\bibitem{groll2021euro2020}
Groll, A., Hvattum, L. M., Ley, C., Popp, F., Schauberger, G.,
Van Eetvelde, H., and Zeileis, A.
\emph{Hybrid Machine Learning Forecasts for the UEFA EURO 2020}.
arXiv:2106.05799, 2021.
\url{https://arxiv.org/abs/2106.05799}

\bibitem{groll2024euro2024}
Groll, A., Hvattum, L. M., Ley, C., Sternemann, J., Schauberger, G.,
and Zeileis, A.
\emph{Modeling and Prediction of the UEFA EURO 2024 via Combined Statistical
Learning Approaches}. arXiv:2410.09068, 2024.
\url{https://arxiv.org/abs/2410.09068}

\bibitem{petretta2021}
Petretta, M., Schiavon, L., and Diquigiovanni, J. \emph{On the Dependence in
Football Match Outcomes: Traditional Model Assumptions and an Alternative
Proposal}. arXiv:2103.07272, 2021.
\url{https://arxiv.org/abs/2103.07272}

\bibitem{michels2023}
Michels, R., {\"O}tting, M., and Karlis, D. \emph{Extending the Dixon and Coles
Model: An Application to Women's Football Data}. arXiv:2307.02139, 2023.
\url{https://arxiv.org/abs/2307.02139}

\bibitem{targets}
Landau, W. M. \emph{The targets R package: a dynamic Make-like function-oriented
pipeline toolkit for reproducibility and high-performance computing}.
\url{https://docs.ropensci.org/targets/}

\bibitem{mass}
Venables, W. N. and Ripley, B. D. \emph{Modern Applied Statistics with S}.
Springer, 2002.
\end{thebibliography}

\end{document}
